\subsection{DEFINITION OF AN ALGEBRA BY GENERATORS AND RELATIONS}

Let \(\mathbf{F}\) be an algebra over \(\mathbf{A}\) and \((x_{i})_{i\in I}\) a family of elements of \(\mathbf{F}\) . By no. 7, Proposition 7, there exists a unique homomorphism \(f\colon \operatorname{Lib}_{\mathbf{A}}(\mathbf{I})\to \mathbf{F}\) such that \(f(\mathbf{X}_i) = x_i\) for all \(i\in I\) . For \(f\) to be surjective, it is necessary and sufficient that \((x_{i})_{i\in I}\) be a generating system of \(\mathbf{F}\) .

If \(\mathbf{U} \in \operatorname{Lib}_{\mathbf{A}}(\mathbf{I})\) , \(f(\mathbf{U})\) is called the element of \(\mathbf{F}\) derived from \(\mathbf{U}\) by substituting the elements \(x_{i}\) for the indeterminates \(\mathbf{X}_{i}\) , or also the value of \(\mathbf{U}\) for the values \(x_{i}\) of the indeterminates \(\mathbf{X}_{i}\) ; it is usually denoted by \(\mathbf{U}((x_{i})_{i \in I})\) ; in particular \(\mathbf{U}((\mathbf{X}_{i})_{i \in I}) = \mathbf{U}\) . If \(\lambda\) is a homomorphism of \(\mathbf{F}\) into an algebra \(\mathbf{F}'\) over \(\mathbf{A}\) , then

\[
\lambda (\mathrm{U} ((x _ {i}) _ {i \in \mathrm{I}})) = \mathrm{U} ((\lambda (x _ {i})) _ {i \in \mathrm{I}}).
\]

Consider in particular the case where \(\mathrm{F} = \operatorname{Lib}_{\mathbf{A}}(\mathbf{J})\) , where J is another set; for every family \((\mathrm{H}_{i})_{i \in I}\) of elements of \(\operatorname{Lib}_{\mathbf{A}}(\mathbf{J})\) and every family \((y_{j}')_{j \in J}\) of elements of an A-algebra \(F'\) ,

\[
\mathrm{U} ((\mathrm{H} _ {i}) _ {i \in I})) ((y _ {j} ^ {\prime}) _ {j \in J}) = \mathrm{U} ((\mathrm{H} _ {i} ((y _ {j} ^ {\prime}) _ {j \in J})) _ {i \in I}).\tag{37}
\]

In the above notation, every element U of \(\operatorname{Lib}_{\mathbf{A}}(\mathbf{I})\) such that \(\mathrm{U}((x_{i})_{i\in\mathbf{I}})=0\) , or also such that \(f(\mathrm{U})=0\) , is called a relator of the family \((x_{i})_{i\in\mathbf{I}}\) in F. The two-sided ideal \(\operatorname{Ker}(f)\) consisting of these elements is called the ideal of relators of \((x_{i})\) .

Let \((\mathbf{R}_j)_{j\in J}\) be a family of elements of \(\operatorname{Lib}_{\mathbf{A}}(\mathbf{I})\) . \(((x_i)_{i\in I}, (\mathbf{R}_j)_{j\in J})\) is called a presentation of the algebra \(F\) if \((x_i)_{i\in I}\) is a generating system of \(F\) and the two-sided ideal of \(\operatorname{Lib}_{\mathbf{A}}(\mathbf{I})\) generated by the \(\mathbf{R}_j\) is equal to the ideal of relators of the family \((x_i)_{i\in I}\) ; the \(x_i\) are called the generators and the \(\mathbf{R}_j\) the relators of the presentation.

Consider now any set I and a family \((\mathbf{R}_j)_{j\in J}\) of elements of \(\operatorname{Lib}_{\mathbf{A}}(\mathbf{I})\) . The quotient algebra E of \(\operatorname{Lib}_{\mathbf{A}}(\mathbf{I})\) by the two-sided ideal generated by the family \((\mathbf{R}_j)\) is called the universal algebra defined by the generating system I related by the family of relators \((\mathbf{R}_j)_{j\in J}\) . Clearly, if \(\overline{\mathbf{X}}_i\) denotes the image of \(\mathbf{X}_i\) in E,

\[
\left(\left(\overline {{\mathbf {X}}} _ {i}\right) _ {i \in I}, \left(\mathbf {R} _ {j}\right) _ {j \in J}\right)
\]

is a presentation of E. Moreover, if \((x_{i})_{i\in I}\) is a family of elements of an algebra F and \(\mathbf{R}_j((x_i)_{i\in I}) = 0\) for all \(j\in J\) , there exists a unique homomorphism \(g:\mathbf{E}\to \mathbf{F}\) such that \(g(\overline{\mathbf{X}}_i) = x_i\) for all \(i\in I\) ; for \(((x_i)_{i\in I},(\mathbf{R}_j)_{j\in J})\) to be a presentation of F, it is necessary and sufficient that \(g\) be bijective.

These remarks justify the following abuse of language: instead of saying `` \(((x_i)_{i \in I}, (R_j)_{j \in J})\) is a presentation of F'', it is also said that ``F is the algebra generated by the generators \(x_i\) subject to the relations \(R_j((x_i)_{i \in I}) = 0\) ''. When the \(R_j\) are of the form \(P_j - Q_j\) , it is also said that ``F is the algebra generated by the \(x_i\) subject to the relations \(P_j((x_i)) = Q_j((x_i))\) ''.

Let H be a set; we shall say that an element S of \(\operatorname{Lib}_{\mathbf{A}}(\mathbf{H})\) is a universal relator for an A-algebra F is \(\mathrm{S}((x_{h})_{h\in\mathbb{H}})=0\) for every family \((x_{h})_{h\in\mathbb{H}}\) of elements of F with H as indexing set.

Examples. (1) Take \(\mathbf{H} = \{1,2,3\}\) ; the algebras which admit

\[
(\mathrm{X} _ {1} \mathrm{X} _ {2}) \mathrm{X} _ {3} - \mathrm{X} _ {1} (\mathrm{X} _ {2} \mathrm{X} _ {3})
\]

as universal relator are the associative algebras. The algebras which admit \(X_{1}X_{2}-X_{2}X_{1}\) as universal relator are the commutative algebras. The algebras which admit the universal relators \(X_{1}X_{1}\) and

\[
(\mathrm{X} _ {1} \mathrm{X} _ {2}) \mathrm{X} _ {3} + (\mathrm{X} _ {2} \mathrm{X} _ {3}) \mathrm{X} _ {1} + (\mathrm{X} _ {3} \mathrm{X} _ {1}) \mathrm{X} _ {2}
\]

as universal relator are the Lie algebras.

Let I be a set; let a family \((\mathrm{S}_{k})_{k\in\mathbb{K}}\) of elements of \(\operatorname{Lib}_{\mathbf{A}}(\mathbf{H})\) be given and consider the set T of elements of \(\operatorname{Lib}_{\mathbf{A}}(\mathbf{I})\) of the form \(\mathrm{S}_{k}((\mathrm{U}_{h}))_{h\in\mathbb{H}}\) , where k runs through K and, for each \(k,\left(\mathrm{U}_{h}\right)_{h\in\mathbb{H}}\) runs through the set of families of elements of \(\operatorname{Lib}_{\mathbf{A}}(\mathbf{I})\) with H as indexing set; consider a family \((\mathrm{R}_{j})_{j\in\mathbb{J}}\) with T as set of elements. Then let F be the universal algebra defined by the generating system I related by the family \((\mathrm{R}_{j})_{j\in\mathbb{J}}\) and let \(u:\operatorname{Lib}_{\mathbf{A}}(\mathbf{I})\to\mathbf{F}\) be the canonical homomorphism, so that \(\operatorname{Ker}(u)\) is generated by the elements \(\mathrm{S}_{k}((\mathrm{U}_{h})_{h\in\mathbb{H}})\) for all \(k\in K\) and all families \((\mathrm{U}_{h})_{h\in\mathbb{H}}\) of elements of \(\operatorname{Lib}_{\mathbf{A}}(\mathbf{I})\) ; clearly each of the \(S_{k}\) \((k\in\mathbb{K})\) is a universal relator for F. Now let \(F'\) be an algebra admitting a generating system \((x_{i})_{i\in\mathbb{I}}\) , for which each of the \(S_{k}\) is a universal relator, and let \(u':\operatorname{Lib}_{\mathbf{A}}(\mathbf{I})\to\mathbf{F}'\) be the homomorphism such that \(u'(\mathrm{X}_{i})=x_{i}\) for all \(i\in I\) ; clearly \(\operatorname{Ker}(u)\subset\operatorname{Ker}(u')\) and hence \(u'\) can be written uniquely in the form \(u'=h\circ u\) , where \(h:F\to F'\) is a homomorphism such that \(h(\overline{\mathrm{X}}_{i})=x_{i}\) for all \(i\in I\) . For this reason F is called the universal algebra defined by the generating system I, corresponding to the family of universal relators \((\mathrm{S}_{k})_{k\in\mathbb{K}}\) . By an abuse of language, F is sometimes called the universal algebra generated by I and subject to the identities \(\mathrm{S}_{k}((u_{h}))=0\) for every family \((u_{h})_{h\in\mathbb{H}}\) of elements of F.

Example 2. Let \(\mathbf{L}'\) be the universal algebra generated by \(\mathbf{I}\) and subject to the identities \((uv)w - u(vw) = 0\) for every family of three elements of \(\mathbf{L}'\) and let \(\mathbf{L}''\) be the algebra obtained by adjoining a unit element to \(\mathbf{L}'\) ; then there exists a unique unital isomorphism \(g\) of \(\mathbf{L}''\) onto \(\operatorname{Libas}_{\mathbf{A}}(\mathbf{I})\) such that \(g(\overline{\mathbf{X}}_i) = \mathbf{X}_i\) for all \(i \in \mathbf{I}\) . For clearly \(\mathbf{L}''\) is associative and the existence of the homomorphism \(g\) follows from the definition of \(\mathbf{L}'\) and the remarks preceding it; but then clearly \(\mathbf{L}''\) satisfies the universal property (no. 7, Proposition 7) which characterizes \(\operatorname{Libas}_{\mathbf{A}}(\mathbf{I})\) , whence the conclusion.

Considerations analogous to the above can be applied to associative algebras (resp. commutative associative algebras), taking account of the following remarks. When the context gives sufficient indication that the algebras considered are unital algebras, by an abuse of language, a family of elements \((x_{i})_{i\in I}\) of an algebra \(F\) , such that the subalgebra generated by the \(x_{i}\) ( \(i\in I\) ) and the unit element is equal to \(F\) , is often called a generating system of \(F\) . Then let \(F\) be a unital associative algebra over \(A\) and let \((x_{i})_{i\in I}\) be a family of elements of \(F\) ; by virtue of no. 7, Proposition 7, there exists a unique unital homomorphism \(f:\mathrm{Libas}_{\mathbf{A}}(\mathbf{I})\to F\) such that \(f(\mathbf{X}_i) = x_i\) for all \(i\in I\) ; if \(\mathbf{U}\in \mathrm{Libas}_{\mathbf{A}}(\mathbf{I}), f(\mathbf{U})\) is also called the element of \(F\) derived from \(\mathbf{U}\) by substituting the elements \(x_{i}\) for the indeterminates \(\mathbf{X}_i\) and it is also denoted by \(\mathrm{U}((x_{i})_{i\in I})\) . Then the notions of relator, presentation and universal relator go over immediately to associative algebras; it suffices simply to replace \(\mathrm{Lib}_{\mathbf{A}}(\mathbf{I})\) throughout by \(\mathrm{Libas}_{\mathbf{A}}(\mathbf{I})\) . The universal unital associative algebra defined by the generating system \(I\) related by the family of relators \((\mathbb{R}_j)_{j\in J}\) is the quotient algebra of \(\mathrm{Libas}_{\mathbf{A}}(\mathbf{I})\) by the two-sided ideal generated by the family \((\mathbb{R}_j)\) . The universal unital associative algebra generated by the generating system \(I\) , corresponding to a family of universal relators is defined similarly. We leave to the reader the task of stating the analogous definitions relative to commutative associative algebras with \(\mathrm{Libasc}_{\mathbf{A}}(\mathbf{I})\) substituted for \(\mathrm{Libas}_{\mathbf{A}}(\mathbf{I})\) .

Example 3. Let L' be the universal unital associative algebra generated by I and subject to the identities uv - vu = 0 for every family of two elements of L'. It is seen, as in Example 2, that L' is canonically isomorphic to \(\text{Libasc}_{\mathbf{A}}(\mathbf{I})\) .

